\section{What the Features Know}
\label{sec:features}

\subsection{Encoder comparison}
\label{sec:encoders}

All probes in this section see only images, use the same small probe, and are
trained on all 92k scenes for 100k steps. Only the input changes.
Table~\ref{tab:encoder-results} gives the results, with the no-input and
ego-status probes as reference points.

\begin{table}[htbp]
\centering
\caption{Image encoders under the same probe. Validation: 2000 scenes of the
held-out logs, 16 samples, final checkpoint. navtest: all 12,146 scenes, 32
samples, for the runs we evaluated offline. Energy and
displacements in metres, lower is better.}
\label{tab:encoder-results}
\begin{tblr}{colspec={lrrrrr}, row{1}={font=\bfseries}, cells={font=\small}}
\toprule
Input & Val energy & navtest energy & Val minADE & ADE of mean & Dispersion \\
\midrule
none                             & 13.10 & 13.58 & 1.95 & 5.91 & 0.92 \\
$q_{\text{rec}}$                 & 7.34  &       & 1.01 & 3.26 & 0.94 \\
$q_{\text{sem}}$ + $q_{\text{rec}}$ & 5.62 &     & 0.82 & 2.47 & 0.92 \\
$q_{\text{sem}}$                 & 5.55  & 6.08  & 0.79 & 2.46 & 0.93 \\
DINOv2                           & 5.37  & 5.90  & 0.79 & 2.36 & 0.90 \\
V-JEPA 2                         & 4.66  & 4.88  & 0.68 & 1.99 & 0.90 \\
WM23                             & 3.70  & 4.06  & 0.58 & 1.56 & 0.90 \\
WM11-ctx                         & \textbf{2.91} & \textbf{3.06} & \textbf{0.46} & \textbf{1.19} & 0.88 \\
\midrule
ego status                       & 2.73  &       & 0.45 & 1.01 & 0.90 \\
ego status + command             & 2.39  & 2.48  & 0.40 & 0.90 & 0.91 \\
\bottomrule
\end{tblr}
\end{table}

The Orbis world model is the best image encoder by a wide margin. Block 11 at the
last observed frame lowers the energy score by 37\% against V-JEPA~2 and by 46\%
against DINOv2. The ranking follows how much temporal context the encoder
processes and what it was trained for: per-frame encoders ($q_{\text{sem}}$,
DINOv2) are worst, a general video encoder (V-JEPA~2) is better, and a video model
trained to predict driving is best. The tokenizer latents are much weaker than
the world model built on them, so the world model adds the information itself,
by relating the frames to each other.

Two smaller results explain the gap between the per-frame encoders. The detail
latent $q_{\text{rec}}$ is much worse than $q_{\text{sem}}$ and adds nothing to
it, as expected for a latent trained for pixel reconstruction. And DINOv2's small
lead over $q_{\text{sem}}$ comes from width: reduced to the same 16 channels,
DINOv2 is slightly worse (5.81 against 5.55), which fits the fact that
$q_{\text{sem}}$ is trained to imitate DINOv2. For per-frame encoders, more
history also helps: $q_{\text{sem}}$ of one, two and four frames gives 6.10, 5.83
and 5.55, since speed is only visible as a change between frames.

The dispersion is between 0.88 and 0.94 for every input. At full data the encoder
changes how accurate the distribution is, not how wide it is relative to its
error.

The ego status alone still beats every image encoder on the energy score. This
is the shortcut of Section~\ref{sec:baselines}: speed and acceleration explain
most of the variance of a 4\,s trajectory. Section~\ref{sec:planning} shows that
it does not carry over to driving quality.

\subsection{Where in the world model}
\label{sec:depth}

Figure~\ref{fig:depth} shows the energy score of probes on six blocks of the world
model, for both token positions. Three results stand out.

\begin{figure}[htbp]
\centering
\includegraphics[width=0.72\textwidth]{../pictures/wm_depth.png}
\caption{Validation energy of the probe on world-model features against block
index, for the next-frame slot and the last observed frame. Dashed lines are the
other encoders and the ego-status probe (without command).}
\label{fig:depth}
\end{figure}

The middle of the network is best. The energy falls steeply up to block 11 and
then rises slowly. The last blocks are closer to the next-frame prediction that
the model must output, and keep less of the driving state.

The last observed frame beats the next-frame slot at every depth, by 0.4 to 0.8 in
energy. We expected the opposite, since the next-frame slot is where the model
forms its guess of the future. But at $\tau=1$ that slot starts as pure noise, and
after one forward pass it holds a mix of noise and prediction. The model keeps
its estimate of the motion in the context it reads from.

Every block from 7 on beats every other encoder, and even block 3 beats
$q_{\text{sem}}$ and DINOv2. This rules out the possibility that one lucky layer
explains the result. The whole world model holds more about the future
trajectory than the image encoders.

\subsection{Linear read-out}
\label{sec:linear}

The probe is nonlinear and generative. To see what is directly available, we
fitted ridge regressions from the features (pooled to $3\times4$ cells, so every
input has 768 numbers) to quantities of the current scene and its future, and
scored them on the validation logs (Table~\ref{tab:linear},
Figure~\ref{fig:linear}).

\begin{table}[htbp]
\centering
\caption{Linear read-out from 768 feature numbers, $R^2$ on validation logs.
Command is the 3-class accuracy (majority class 0.63). Progress, lateral offset
and heading are those of the pose at 4\,s.}
\label{tab:linear}
\begin{tblr}{colspec={lrrrrrr}, row{1}={font=\bfseries}}
\toprule
Feature & Speed & Accel. & Progress & Lateral & Heading & Command \\
\midrule
$q_{\text{rec}}$  & 0.36 & 0.24 & 0.32 & 0.41 & 0.38 & 0.86 \\
$q_{\text{sem}}$  & 0.59 & 0.40 & 0.56 & 0.48 & 0.45 & 0.89 \\
DINOv2, 16 ch.    & 0.57 & 0.39 & 0.54 & 0.48 & 0.45 & 0.89 \\
V-JEPA 2          & 0.87 & 0.48 & 0.74 & 0.75 & 0.66 & 0.91 \\
WM11              & 0.90 & 0.57 & 0.80 & 0.75 & 0.67 & 0.91 \\
WM11-ctx          & \textbf{0.95} & \textbf{0.68} & \textbf{0.87} & \textbf{0.84} & \textbf{0.74} & 0.91 \\
WM23-ctx          & 0.94 & 0.50 & 0.84 & 0.81 & 0.72 & 0.92 \\
\bottomrule
\end{tblr}
\end{table}

\begin{figure}[htbp]
\centering
\includegraphics[width=\textwidth]{../pictures/linear_probe_depth.png}
\caption{Linear read-out $R^2$ against world-model block, for the two token
positions. Dashed lines are the other encoders.}
\label{fig:linear}
\end{figure}

The linear results agree with the probe. The world model holds the ego speed
almost perfectly ($R^2=0.95$), against 0.87 for V-JEPA~2 and 0.59 for
$q_{\text{sem}}$. Acceleration is the hardest quantity and peaks sharply at
block 11 (0.68), falling to 0.50 at the last block. So the middle of the network
holds the most dynamics. The last observed frame is the better position for every
quantity at every depth. The driving command is the only target where all
encoders are close (0.86 to 0.92), because the route is mostly not visible in the
front camera. This is also why giving the probe the command helps
(Section~\ref{sec:planning}).

\subsection{Generalization}
\label{sec:generalization}

\paragraph{Held-out logs against navtest.}
All model selection used held-out training logs. The navtest energy is 3 to 10\%
higher than the validation energy for every probe, and the ranking of the probes
is the same on both (Table~\ref{tab:encoder-results}). So the log-level split is a
reliable stand-in for the test set, and the conclusions above are not tuned to
the validation data.

\paragraph{Unseen city.}
A $q_{\text{sem}}$ probe trained without Singapore, the only left-hand-traffic
city, is 50\% worse there than the probe trained on all cities (energy 9.43
against 6.31). On the US cities the two are equal. The dispersion in Singapore
falls from 0.86 to 0.66. The model does not widen its distribution on the
unfamiliar city. It samples with the usual width around the wrong trajectory.

\paragraph{Unseen combinations.}
A harder test is whether the probe combines factors it has seen separately. We
removed two (city, command) pairs from training: left turns in Singapore and right
turns in Boston. The model still sees left turns in the US cities and Singapore
scenes going straight or right. Both models get the command as an input.

\begin{table}[htbp]
\centering
\caption{Compositional test, $q_{\text{sem}}$ + command probe, 32 samples on
validation scenes of the held-out pairs (314 Singapore-left, 260 Boston-right).
The reference model saw these pairs in training.}
\label{tab:compositional}
\begin{tblr}{colspec={llrrr}, row{1}={font=\bfseries}}
\toprule
Pair & Model & Energy & minADE & Spread \\
\midrule
Singapore, left & reference & 4.51 & 0.57 & 2.67 \\
Singapore, left & held out  & 9.11 & 1.39 & 2.58 \\
Boston, right   & reference & 6.40 & 0.91 & 2.57 \\
Boston, right   & held out  & 7.57 & 1.11 & 2.50 \\
\bottomrule
\end{tblr}
\end{table}

On the rest of the validation set the held-out model is only 5\% worse, about
what the removed data costs. On Singapore left turns its energy doubles
(Table~\ref{tab:compositional}). A left turn in Singapore is the short near-side
turn, which has the geometry of a right turn in the US, so knowing US left turns
does not help. Boston right turns look like right turns in the other US cities
and lose only 18\%. The spread of the samples is unchanged, so the model gives no
sign that the input is new to it. Longer training makes this worse: the energy
on the held-out pairs is best at 40k steps (7.64) and rises to 8.41 at 100k,
while the reference model keeps improving. The probe fits the combinations it
has seen more and more tightly, the same trend as the collapse in
Section~\ref{sec:collapse}.
